\chapter{Venous Insufficiency}

\section*{Definition}
Chronic venous insufficiency (CVI) is a condition of impaired venous return from the lower extremities due to valvular incompetence, venous obstruction, or calf muscle pump dysfunction. It encompasses the clinical spectrum from varicose veins and edema to lipodermatosclerosis and venous ulceration.

\section*{Pathophysiology}
Failed venous valves allow retrograde flow (reflux) of blood, producing sustained venous hypertension in the lower extremity during standing. The resultant capillary hypertension causes fluid extravasation, hemosiderin deposition, activation of matrix metalloproteinases, and ultimately tissue fibrosis and ulceration.

\section*{Etiology}
\begin{itemize}
  \item \textbf{Primary valvular incompetence:} Idopathic, familial predisposition, age-related valve degeneration
  \item \textbf{Secondary (post-thrombotic):} Deep vein thrombosis causing valve destruction and venous scarring (most common cause of severe CVI)
  \item \textbf{Obstructive:} Iliac vein compression (May-Thurner syndrome), extrinsic compression (tumor, pregnancy), venous stenosis
  \item \textbf{Calf muscle pump dysfunction:} Immobility, neuromuscular disorders, arthritis limiting ankle range of motion, obesity
  \item \textbf{Other:} Congenital vein wall weakness (Ehlers-Danlos), arteriovenous fistula, high-output states
\end{itemize}

\section*{Indications for Evaluation}
Seek care if:
\begin{itemize}
  \item Persistent ankle or leg edema not relieved by overnight elevation
  \item Skin changes: red-brown hemosiderin pigmentation, thickening/fibrosis (lipodermatosclerosis), eczematous dermatitis
  \item Atrophie blanche (white, atrophic patches with telangiectasias) or healed/active venous ulcer
  \item Worsening symptoms despite conservative therapy (compression, elevation)
  \item Suspected deep vein obstruction or acute-on-chronic worsening
\end{itemize}

\section*{Related Symptoms}
\begin{itemize}
  \item Leg heaviness, fatigue, aching, edema, itching, burning, skin discoloration, venous ulcers, varicose veins
\end{itemize}
\textbf{Note:} Venous ulcers account for 70--80\% of all lower extremity ulcers, typically occurring in the gaiter area (medial malleolus) and often recurring without adequate compression therapy.

\section*{Self-Care / Home Management}
\begin{itemize}
  \item Daily graduated compression stockings (20--40 mmHg); put on before rising and remove at bedtime
  \item Leg elevation above heart level for 30 minutes, 3--4 times per day, and overnight with bed blocks
  \item Regular ankle range-of-motion exercises and walking to activate the calf muscle pump
  \item Weight reduction and avoidance of prolonged standing or sitting
  \item Meticulous skin care: emollient moisturizers, avoid scratching, prompt treatment of minor injuries
\end{itemize}

\section*{Diagnostic Approach}
\begin{itemize}
  \item History and physical exam with CEAP classification of venous disease
  \item Duplex venous ultrasound with reflux testing: measure venous refill time, reflux duration (greater than 0.5 s for superficial, greater than 1.0 s for deep veins)
  \item Venous plethysmography (air or strain gauge) for global assessment of venous function
  \item CT or MR venography for suspected iliac vein obstruction or pre-procedural planning
  \item Ankle-brachial index (ABI) to exclude concurrent arterial disease before prescribing compression therapy
\end{itemize}

\section*{Treatment / Management Overview}
\begin{itemize}
  \item Compression therapy (graduated stockings, multilayered bandaging) as cornerstone of management
  \item Endovenous ablation (laser, radiofrequency, mechanochemical) for incompetent saphenous veins
  \item Stenting for iliofemoral vein obstruction (post-thrombotic or compressive)
  \item Wound care for venous ulcers: debridement, moist wound healing, compression, and skin grafting for refractory cases
  \item Pentoxifylline or micronized purified flavonoid fraction (MPFF) as adjunctive pharmacotherapy to reduce ulcer healing time
\end{itemize}

\clearpage
